% Options for packages loaded elsewhere
\PassOptionsToPackage{unicode}{hyperref}
\PassOptionsToPackage{hyphens}{url}
%
\documentclass[
]{article}
\usepackage{amsmath,amssymb}
\usepackage{iftex}
\ifPDFTeX
  \usepackage[T1]{fontenc}
  \usepackage[utf8]{inputenc}
  \usepackage{textcomp} % provide euro and other symbols
\else % if luatex or xetex
  \usepackage{unicode-math} % this also loads fontspec
  \defaultfontfeatures{Scale=MatchLowercase}
  \defaultfontfeatures[\rmfamily]{Ligatures=TeX,Scale=1}
\fi
\usepackage{lmodern}
\ifPDFTeX\else
  % xetex/luatex font selection
\fi
% Use upquote if available, for straight quotes in verbatim environments
\IfFileExists{upquote.sty}{\usepackage{upquote}}{}
\IfFileExists{microtype.sty}{% use microtype if available
  \usepackage[]{microtype}
  \UseMicrotypeSet[protrusion]{basicmath} % disable protrusion for tt fonts
}{}
\makeatletter
\@ifundefined{KOMAClassName}{% if non-KOMA class
  \IfFileExists{parskip.sty}{%
    \usepackage{parskip}
  }{% else
    \setlength{\parindent}{0pt}
    \setlength{\parskip}{6pt plus 2pt minus 1pt}}
}{% if KOMA class
  \KOMAoptions{parskip=half}}
\makeatother
\usepackage{xcolor}
\usepackage[margin=1in]{geometry}
\usepackage{color}
\usepackage{fancyvrb}
\newcommand{\VerbBar}{|}
\newcommand{\VERB}{\Verb[commandchars=\\\{\}]}
\DefineVerbatimEnvironment{Highlighting}{Verbatim}{commandchars=\\\{\}}
% Add ',fontsize=\small' for more characters per line
\usepackage{framed}
\definecolor{shadecolor}{RGB}{248,248,248}
\newenvironment{Shaded}{\begin{snugshade}}{\end{snugshade}}
\newcommand{\AlertTok}[1]{\textcolor[rgb]{0.94,0.16,0.16}{#1}}
\newcommand{\AnnotationTok}[1]{\textcolor[rgb]{0.56,0.35,0.01}{\textbf{\textit{#1}}}}
\newcommand{\AttributeTok}[1]{\textcolor[rgb]{0.13,0.29,0.53}{#1}}
\newcommand{\BaseNTok}[1]{\textcolor[rgb]{0.00,0.00,0.81}{#1}}
\newcommand{\BuiltInTok}[1]{#1}
\newcommand{\CharTok}[1]{\textcolor[rgb]{0.31,0.60,0.02}{#1}}
\newcommand{\CommentTok}[1]{\textcolor[rgb]{0.56,0.35,0.01}{\textit{#1}}}
\newcommand{\CommentVarTok}[1]{\textcolor[rgb]{0.56,0.35,0.01}{\textbf{\textit{#1}}}}
\newcommand{\ConstantTok}[1]{\textcolor[rgb]{0.56,0.35,0.01}{#1}}
\newcommand{\ControlFlowTok}[1]{\textcolor[rgb]{0.13,0.29,0.53}{\textbf{#1}}}
\newcommand{\DataTypeTok}[1]{\textcolor[rgb]{0.13,0.29,0.53}{#1}}
\newcommand{\DecValTok}[1]{\textcolor[rgb]{0.00,0.00,0.81}{#1}}
\newcommand{\DocumentationTok}[1]{\textcolor[rgb]{0.56,0.35,0.01}{\textbf{\textit{#1}}}}
\newcommand{\ErrorTok}[1]{\textcolor[rgb]{0.64,0.00,0.00}{\textbf{#1}}}
\newcommand{\ExtensionTok}[1]{#1}
\newcommand{\FloatTok}[1]{\textcolor[rgb]{0.00,0.00,0.81}{#1}}
\newcommand{\FunctionTok}[1]{\textcolor[rgb]{0.13,0.29,0.53}{\textbf{#1}}}
\newcommand{\ImportTok}[1]{#1}
\newcommand{\InformationTok}[1]{\textcolor[rgb]{0.56,0.35,0.01}{\textbf{\textit{#1}}}}
\newcommand{\KeywordTok}[1]{\textcolor[rgb]{0.13,0.29,0.53}{\textbf{#1}}}
\newcommand{\NormalTok}[1]{#1}
\newcommand{\OperatorTok}[1]{\textcolor[rgb]{0.81,0.36,0.00}{\textbf{#1}}}
\newcommand{\OtherTok}[1]{\textcolor[rgb]{0.56,0.35,0.01}{#1}}
\newcommand{\PreprocessorTok}[1]{\textcolor[rgb]{0.56,0.35,0.01}{\textit{#1}}}
\newcommand{\RegionMarkerTok}[1]{#1}
\newcommand{\SpecialCharTok}[1]{\textcolor[rgb]{0.81,0.36,0.00}{\textbf{#1}}}
\newcommand{\SpecialStringTok}[1]{\textcolor[rgb]{0.31,0.60,0.02}{#1}}
\newcommand{\StringTok}[1]{\textcolor[rgb]{0.31,0.60,0.02}{#1}}
\newcommand{\VariableTok}[1]{\textcolor[rgb]{0.00,0.00,0.00}{#1}}
\newcommand{\VerbatimStringTok}[1]{\textcolor[rgb]{0.31,0.60,0.02}{#1}}
\newcommand{\WarningTok}[1]{\textcolor[rgb]{0.56,0.35,0.01}{\textbf{\textit{#1}}}}
\usepackage{longtable,booktabs,array}
\usepackage{calc} % for calculating minipage widths
% Correct order of tables after \paragraph or \subparagraph
\usepackage{etoolbox}
\makeatletter
\patchcmd\longtable{\par}{\if@noskipsec\mbox{}\fi\par}{}{}
\makeatother
% Allow footnotes in longtable head/foot
\IfFileExists{footnotehyper.sty}{\usepackage{footnotehyper}}{\usepackage{footnote}}
\makesavenoteenv{longtable}
\usepackage{graphicx}
\makeatletter
\newsavebox\pandoc@box
\newcommand*\pandocbounded[1]{% scales image to fit in text height/width
  \sbox\pandoc@box{#1}%
  \Gscale@div\@tempa{\textheight}{\dimexpr\ht\pandoc@box+\dp\pandoc@box\relax}%
  \Gscale@div\@tempb{\linewidth}{\wd\pandoc@box}%
  \ifdim\@tempb\p@<\@tempa\p@\let\@tempa\@tempb\fi% select the smaller of both
  \ifdim\@tempa\p@<\p@\scalebox{\@tempa}{\usebox\pandoc@box}%
  \else\usebox{\pandoc@box}%
  \fi%
}
% Set default figure placement to htbp
\def\fps@figure{htbp}
\makeatother
\setlength{\emergencystretch}{3em} % prevent overfull lines
\providecommand{\tightlist}{%
  \setlength{\itemsep}{0pt}\setlength{\parskip}{0pt}}
\setcounter{secnumdepth}{5}
% definitions for citeproc citations
\NewDocumentCommand\citeproctext{}{}
\NewDocumentCommand\citeproc{mm}{%
  \begingroup\def\citeproctext{#2}\cite{#1}\endgroup}
\makeatletter
 % allow citations to break across lines
 \let\@cite@ofmt\@firstofone
 % avoid brackets around text for \cite:
 \def\@biblabel#1{}
 \def\@cite#1#2{{#1\if@tempswa , #2\fi}}
\makeatother
\newlength{\cslhangindent}
\setlength{\cslhangindent}{1.5em}
\newlength{\csllabelwidth}
\setlength{\csllabelwidth}{3em}
\newenvironment{CSLReferences}[2] % #1 hanging-indent, #2 entry-spacing
 {\begin{list}{}{%
  \setlength{\itemindent}{0pt}
  \setlength{\leftmargin}{0pt}
  \setlength{\parsep}{0pt}
  % turn on hanging indent if param 1 is 1
  \ifodd #1
   \setlength{\leftmargin}{\cslhangindent}
   \setlength{\itemindent}{-1\cslhangindent}
  \fi
  % set entry spacing
  \setlength{\itemsep}{#2\baselineskip}}}
 {\end{list}}
\usepackage{calc}
\newcommand{\CSLBlock}[1]{\hfill\break\parbox[t]{\linewidth}{\strut\ignorespaces#1\strut}}
\newcommand{\CSLLeftMargin}[1]{\parbox[t]{\csllabelwidth}{\strut#1\strut}}
\newcommand{\CSLRightInline}[1]{\parbox[t]{\linewidth - \csllabelwidth}{\strut#1\strut}}
\newcommand{\CSLIndent}[1]{\hspace{\cslhangindent}#1}
\usepackage{xcolor}
\definecolor{pending}{RGB}{180,0,0}
\usepackage{bookmark}
\IfFileExists{xurl.sty}{\usepackage{xurl}}{} % add URL line breaks if available
\urlstyle{same}
\hypersetup{
  pdftitle={CyFj11: FlowJo v11 Workspace Import and Legacy Format Export for R-Based Flow Cytometry Analysis},
  hidelinks,
  pdfcreator={LaTeX via pandoc}}

\title{CyFj11: FlowJo v11 Workspace Import and Legacy Format Export for
R-Based Flow Cytometry Analysis}
\author{}
\date{\vspace{-2.5em}}

\begin{document}
\maketitle

\begin{center}\rule{0.5\linewidth}{0.5pt}\end{center}

\section{Title Page}\label{title-page}

\textbf{Title:} CyFj11: FlowJo v11 Workspace Import and Legacy Format
Export for R-Based Flow Cytometry Analysis

\textbf{Short running title:} CyFj11: FlowJo v11 / R Interoperability

\textbf{Manuscript version:} v0.2-draft (2026-04-10)

\textbf{Authors:} {[}TODO: full author list with email addresses{]}

\textbf{Affiliations:} {[}TODO: institutional affiliations{]}

\textbf{Corresponding author:} {[}TODO: name, address, email, ORCID
iD{]}

\textbf{Acknowledgments:} {[}TODO: funding sources, acknowledgments{]}

\textbf{Funding:} {[}TODO: grant numbers and funding bodies{]}

\textbf{Conflict of interest:} The authors declare no conflicts of
interest.

\textbf{Data and code availability:} The CyFj11 package is available at
\textcolor{pending}{[TODO: GitHub/Bioconductor URL]}. All test datasets
and validation scripts are deposited at
\textcolor{pending}{[TODO: Zenodo DOI]}.

\textbf{AI use disclosure:} AI-assisted tools were used for spelling and
grammar checking only. All scientific content and interpretation are the
responsibility of the authors.

\textbf{Author contributions (CRediT):} {[}TODO: assign CRediT roles per
author{]}

\begin{center}\rule{0.5\linewidth}{0.5pt}\end{center}

\section{Abstract}\label{abstract}

Flow cytometry data analysis increasingly requires movement between the
commercial FlowJo platform and the open-source R/Bioconductor ecosystem.
The release of FlowJo version 11 (v11) introduced a JSON-based workspace
format that is incompatible with existing R interoperability tools,
creating a barrier for computational cytometry workflows. We present
CyFj11, an R package that enables full import of FlowJo v11 workspaces
into R \texttt{GatingSet} objects and export to the FlowJo v10 XML
format. Import supports all major gate types (rectangle, polygon,
ellipsoid, range, quadrant, and Boolean) and all common transformation
functions (biexponential, logicle, log, arcsinh, and linear), as well as
compensation matrix extraction. Export targets the FlowJo v10 XML
format, which is Gating-ML 2.0 compliant and supported by current FlowJo
versions, providing a practical interoperability pathway while avoiding
the technical barriers of the undocumented v11 JSON schema. The
asymmetric design of CyFj11 is intentional: v11 import is complete,
while export uses the well-specified v10 standard.
\textcolor{pending}{Validation across 13 structured test scenarios — covering all supported gate types, three nesting depths, and all five transformation functions — demonstrates that R-side population counts are internally consistent across the import/export pipeline. FlowJo-side population count comparison and gate coordinate accuracy measurements are ongoing. [TODO: insert final validation numbers once FlowJo-side counts are recorded.]}
CyFj11 is implemented in pure R, requires no compiled code or container
infrastructure, and runs on all platforms supported by R including macOS
on Apple Silicon. The package is available at
\textcolor{pending}{[TODO: URL]}.

\textbf{Keywords:} flow cytometry, FlowJo, GatingSet, Gating-ML,
workspace interoperability, R/Bioconductor

\begin{center}\rule{0.5\linewidth}{0.5pt}\end{center}

\section{Introduction}\label{introduction}

\subsection{Flow cytometry in modern
biology}\label{flow-cytometry-in-modern-biology}

Flow cytometry is a cornerstone technology in cell biology, immunology,
and clinical diagnostics, enabling simultaneous measurement of multiple
markers across millions of individual cells {[}1{]}. As panel sizes grow
--- from conventional 8--18 parameter fluorescence configurations to
modern spectral cytometry panels exceeding 40 parameters --- the
analytical demands have correspondingly increased. Manual gating in
commercial software remains the dominant practice for initial data
exploration, while computational R-based pipelines are increasingly used
for downstream statistical analysis, batch processing, and reproducible
reporting.

\subsection{FlowJo and the R flowWorkspace
ecosystem}\label{flowjo-and-the-r-flowworkspace-ecosystem}

FlowJo (BD Biosciences) is one of the most widely used commercial
platforms for flow cytometry data analysis, providing a graphical
interface for gating, compensation, and population statistics {[}2{]}.
The R/Bioconductor ecosystem, centered on the \texttt{flowCore} and
\texttt{flowWorkspace} packages, provides a complementary suite of tools
for programmatic analysis: automation of repetitive tasks, integration
with statistical frameworks, and reproducible computational workflows
{[}1,3{]}. Effective interoperability between these two environments ---
enabling a researcher to import a FlowJo gating strategy into R, or to
export R-derived gates back into FlowJo for visual review --- is
therefore essential for many laboratories.

\subsection{The FlowJo v11 interoperability
gap}\label{the-flowjo-v11-interoperability-gap}

The release of FlowJo version 11 introduced a fundamental change in
workspace architecture. The XML-based \texttt{.wsp} format used in v10,
which conforms to the ISAC Gating-ML 2.0 standard {[}4{]}, was replaced
by a JSON-based \texttt{.flowjo} format distributed as a ZIP archive.
This new format uses extensive UUID cross-referencing among population
definitions, sample-specific gate instances, data sources, and group
assignments (see Annex A for a detailed comparison). Because the v11
JSON schema is proprietary and not publicly documented, existing
interoperability tools cannot read or write v11 workspaces.

The primary existing tool for FlowJo--R interoperability,
\texttt{CytoML} (part of the \texttt{flowWorkspace} ecosystem), supports
reading FlowJo v10 workspaces but not v11 {[}5{]}. Write support in
\texttt{CytoML} relies on compiled C++ components distributed via Docker
containers, which have not been maintained for current operating system
versions, including macOS on Apple Silicon (ARM architecture). These
factors create a practical barrier for researchers who work with FlowJo
v11 and need to integrate their analyses with R.

\subsection{What CyFj11 provides}\label{what-cyfj11-provides}

CyFj11 addresses this gap through asymmetric format conversion: full
import support for FlowJo v11 workspaces into R \texttt{GatingSet}
objects, and export to the FlowJo v10 XML format. The asymmetry is
intentional --- v11 import was implemented by reverse-engineering the
ZIP-based JSON structure, while export targets the documented Gating-ML
2.0 v10 standard rather than the opaque v11 schema. FlowJo v11 can open
v10-format workspaces, making this a practical workflow: import v11 for
R analysis, export to v10 for return to FlowJo. The pure R
implementation requires no compiled code or containers, ensuring
cross-platform compatibility.

The tool is positioned as a complement to, rather than a replacement
for, \texttt{CytoML}: it covers the v11 import case that \texttt{CytoML}
does not handle, and provides a maintained native-R v10 export path.

\begin{center}\rule{0.5\linewidth}{0.5pt}\end{center}

\section{The CyFj11 Package}\label{the-cyfj11-package}

\subsection{Design rationale}\label{design-rationale}

CyFj11 provides asymmetric conversion: complete FlowJo v11 import, with
export limited to the FlowJo v10 XML format (Figure 1). This design
reflects a deliberate engineering decision, not an incomplete
implementation. The FlowJo v11 JSON schema requires maintaining a
consistent graph of UUID relationships across dozens of entity types;
without vendor documentation, generating a valid v11 export would
require full reverse-engineering of this schema with no guarantee of
compatibility across FlowJo updates. In contrast, the v10 XML format
follows the published Gating-ML 2.0 standard, enabling confident
generation of valid, interoperable files. A detailed technical
justification is provided in Annex A.

\textcolor{pending}{[TODO: Figure 1 — schematic of the CyFj11 workflow: FlowJo v11 → CyFj11 → GatingSet list (per sample) → merge\_list\_to\_gs() → GatingSet → CyFj11 → FlowJo v10]}

\subsection{Core workflow}\label{core-workflow}

The typical CyFj11 workflow consists of three steps:

\begin{Shaded}
\begin{Highlighting}[]
\FunctionTok{library}\NormalTok{(CyFj11)}

\CommentTok{\# Step 1: Read a FlowJo v11 workspace}
\NormalTok{ws }\OtherTok{\textless{}{-}} \FunctionTok{read\_flowjo11\_workspace}\NormalTok{(}\StringTok{"experiment.flowjo"}\NormalTok{)}

\CommentTok{\# Step 2: Convert to a GatingSet (resolves FCS files, applies compensations}
\CommentTok{\#         and transformations, executes gates)}
\NormalTok{gs }\OtherTok{\textless{}{-}} \FunctionTok{fj11\_to\_gatingset}\NormalTok{(ws, }\AttributeTok{group\_name =} \StringTok{"All Samples"}\NormalTok{)}

\CommentTok{\# Step 3: Export back to FlowJo v10 format}
\FunctionTok{export\_flowjo10\_workspace}\NormalTok{(gs, }\AttributeTok{output\_path =} \StringTok{"."}\NormalTok{, }\AttributeTok{workspace\_name =} \StringTok{"exported"}\NormalTok{)}
\end{Highlighting}
\end{Shaded}

\texttt{read\_flowjo11\_workspace()} parses the ZIP archive and extracts
the JSON workspace structure. \texttt{fj11\_to\_gatingset()} locates FCS
files, applies the compensation matrix, maps transformations, and
constructs a named list of per-sample \texttt{GatingSet} objects,
preserving any sample-specific gate variations from FlowJo v11's
tailored gate mechanism. Because certain properties --- including
per-sample gate adjustments and transformation parameters --- can differ
between samples, the result is a list rather than a single merged
\texttt{GatingSet}. Users can combine samples for downstream analysis
using \texttt{flowWorkspace::merge\_list\_to\_gs()}, with the
understanding that this step may discard per-sample variation.
\texttt{export\_flowjo10\_workspace()} generates Gating-ML
2.0--compliant XML that FlowJo v11 can open.

\subsection{Supported features}\label{supported-features}

Table 1 summarizes gate types and transformation functions supported for
import and export. All types supported on import can be inspected and
reanalyzed in R; the export column indicates whether that type can be
written back to a v10 workspace.

\textbf{Table 1. Gate types and transformations supported by CyFj11.}

\begin{longtable}[]{@{}lll@{}}
\toprule\noalign{}
Feature & Import (v11 → R) & Export (R → v10) \\
\midrule\noalign{}
\endhead
\bottomrule\noalign{}
\endlastfoot
\textbf{Gate types} & & \\
Rectangle (1D and 2D) & Yes & Yes \\
Polygon & Yes & Yes \\
Ellipsoid & Yes & Yes \\
Range (1D) & Yes & No\textsuperscript{a} \\
Quadrant & Yes & No\textsuperscript{b} \\
Boolean (AND/OR/NOT) & Yes & Yes \\
\textbf{Transformations} & & \\
Biexponential (biexp) & Yes & Yes \\
Logicle & Yes & Yes \\
Log & Yes & Yes \\
Arcsinh (fasinh) & Yes & Yes \\
Linear & Yes & Yes \\
\textbf{Other} & & \\
Compensation matrix & Yes & Yes \\
Per-sample (tailored) gates & Yes & Merged to
group\textsuperscript{c} \\
Sample keywords / metadata & Yes & Yes \\
\end{longtable}

\textsuperscript{a} Range gates are converted to 1D rectangle gates on
export. \textsuperscript{b} Quadrant gates have no direct v10 equivalent
and are excluded from export with a warning. \textsuperscript{c}
Per-sample gate variations are not expressible in v10 format; the
group-level gate is used.

\subsection{Additional functionality}\label{additional-functionality}

\texttt{pretty\_print\_flowjo()} extracts and formats the JSON content
of a \texttt{.flowjo} file to a human-readable text file, which is
useful for debugging and workspace inspection.

Verbose logging can be enabled with \texttt{set\_verbose(TRUE)} to trace
conversion steps during troubleshooting; the current verbosity state can
be queried with \texttt{get\_verbose()}.

\subsection{Installation}\label{installation}

CyFj11 requires R ≥ 4.0 and Bioconductor packages \texttt{flowCore} and
\texttt{flowWorkspace}. It can be installed from
\textcolor{pending}{[TODO: GitHub/Bioconductor URL]}:

\begin{Shaded}
\begin{Highlighting}[]
\CommentTok{\# [}\AlertTok{TODO}\CommentTok{: install instructions — GitHub / Bioconductor]}
\end{Highlighting}
\end{Shaded}

\begin{center}\rule{0.5\linewidth}{0.5pt}\end{center}

\section{Validation}\label{validation}

\subsection{Overview}\label{overview}

We validated CyFj11 using 13 structured test scenarios designed to
exercise all supported gate types and transformations in isolation and
in combination (Table 2; full descriptions in Annex B). For each
scenario, synthetic FCS files were generated, a FlowJo v11 workspace was
constructed programmatically, converted to a \texttt{GatingSet} via
\texttt{fj11\_to\_gatingset()}, and exported to v10 format via
\texttt{export\_flowjo10\_workspace()}. Population counts from R were
taken as the ground truth for each scenario. FlowJo-side validation
(opening the exported v10 workspace in FlowJo and comparing reported
counts) is described below.

Test workspaces, FCS files, and analysis scripts are deposited at
\textcolor{pending}{[TODO: Zenodo DOI]}.

\subsection{Population count preservation --- R-side
baseline}\label{population-count-preservation-r-side-baseline}

Table 2 reports the population counts produced by R for each test
scenario. These serve as the expected values for downstream FlowJo
comparison.

\textbf{Table 2. Population counts from R GatingSet across 13 validation
scenarios.}

\begin{longtable}[]{@{}
  >{\raggedright\arraybackslash}p{(\linewidth - 8\tabcolsep) * \real{0.2000}}
  >{\raggedright\arraybackslash}p{(\linewidth - 8\tabcolsep) * \real{0.2000}}
  >{\raggedright\arraybackslash}p{(\linewidth - 8\tabcolsep) * \real{0.2000}}
  >{\raggedright\arraybackslash}p{(\linewidth - 8\tabcolsep) * \real{0.2000}}
  >{\raggedright\arraybackslash}p{(\linewidth - 8\tabcolsep) * \real{0.2000}}@{}}
\toprule\noalign{}
\begin{minipage}[b]{\linewidth}\raggedright
Test
\end{minipage} & \begin{minipage}[b]{\linewidth}\raggedright
Gate type
\end{minipage} & \begin{minipage}[b]{\linewidth}\raggedright
Transform
\end{minipage} & \begin{minipage}[b]{\linewidth}\raggedright
Population
\end{minipage} & \begin{minipage}[b]{\linewidth}\raggedright
R count
\end{minipage} \\
\midrule\noalign{}
\endhead
\bottomrule\noalign{}
\endlastfoot
01 & None (root only) & Linear & root & 5,000 \\
02 & 1D rectangle & Linear & FSC\_filter & 4,776 \\
03 & 2D rectangle & Linear & cells & 4,631 \\
04 & Polygon & Linear & singlets & 7,555 \\
05 & Ellipsoid & Linear & ellipse\_cells & 7,301 \\
06 & 1D rectangle & Biexponential & FITC\_pos & 2,403 \\
07 & 1D rectangle & Log & PE\_pos & 3,198 \\
08 & 1D rectangle & Arcsinh & APC\_pos & 1,606 \\
09 & Hierarchical (2 levels) & Linear & cells / singlets & 9,441 /
8,800 \\
10 & Boolean AND & Linear & FSC\_gate / SSC\_gate / both & 7,554 / 7,392
/ 6,979 \\
11 & Mixed hierarchy (biexp + arcsinh + boolean) & Mixed & cells /
FITC\_pos / APC\_pos / double\_pos & 9,696 / 2,897 / 1,930 / 1,928 \\
12 & Ellipsoid + range + AND/OR/NOT & Log + biexp & cells / FITC\_hi /
PE\_hi / double\_hi / either\_hi / neither & 9,765 / 2,934 / 3,831 /
2,879 / 3,886 / 5,879 \\
13 & 3-level hierarchy (polygon/rect/ellipse/boolean) & Arcsinh & live /
singlets / FITC\_pos / double\_pos / not\_double & 9,939 / 9,470 / 2,838
/ 1,887 / 7,583 \\
\end{longtable}

\subsection{Population count preservation --- FlowJo
comparison}\label{population-count-preservation-flowjo-comparison}

\textcolor{pending}{[TODO: Open each \texttt{export.wsp} in FlowJo v11 (version [TODO]) and record the displayed population counts. Add a column "FlowJo count" and "Difference" to Table 2. Expected: differences ≤ 0.5\% for all populations except those in Test 12 (log transform), where coordinate drift has been documented (see Limitations).]}

\subsection{Gate coordinate accuracy}\label{gate-coordinate-accuracy}

Gate boundary coordinates were extracted from each exported v10 XML file
and compared to the original gate definitions used to construct the test
workspaces (Table 3). For linear rectangle and polygon gates, XML
coordinates (in raw data space) were compared directly to the original
flowCore gate coordinates. For transformed gates (biexponential, log,
arcsinh), the XML data-space coordinates were re-transformed using the
same forward transform function and compared to the original
transformed-space gate boundaries. Validation was performed using the
\texttt{validate\_gate\_coordinates.R} script (included in
\texttt{inst/scripts/};
\textcolor{pending}{[TODO: confirm final installation path]}).

\textbf{Table 3. Gate coordinate accuracy across 8 test scenarios
covering 5 gate/transform combinations.}

\begin{longtable}[]{@{}
  >{\raggedright\arraybackslash}p{(\linewidth - 14\tabcolsep) * \real{0.1250}}
  >{\raggedright\arraybackslash}p{(\linewidth - 14\tabcolsep) * \real{0.1250}}
  >{\raggedright\arraybackslash}p{(\linewidth - 14\tabcolsep) * \real{0.1250}}
  >{\raggedright\arraybackslash}p{(\linewidth - 14\tabcolsep) * \real{0.1250}}
  >{\raggedright\arraybackslash}p{(\linewidth - 14\tabcolsep) * \real{0.1250}}
  >{\raggedright\arraybackslash}p{(\linewidth - 14\tabcolsep) * \real{0.1250}}
  >{\raggedright\arraybackslash}p{(\linewidth - 14\tabcolsep) * \real{0.1250}}
  >{\raggedright\arraybackslash}p{(\linewidth - 14\tabcolsep) * \real{0.1250}}@{}}
\toprule\noalign{}
\begin{minipage}[b]{\linewidth}\raggedright
Test
\end{minipage} & \begin{minipage}[b]{\linewidth}\raggedright
Population
\end{minipage} & \begin{minipage}[b]{\linewidth}\raggedright
Gate type
\end{minipage} & \begin{minipage}[b]{\linewidth}\raggedright
Transform
\end{minipage} & \begin{minipage}[b]{\linewidth}\raggedright
N coordinates
\end{minipage} & \begin{minipage}[b]{\linewidth}\raggedright
Mean deviation
\end{minipage} & \begin{minipage}[b]{\linewidth}\raggedright
Max deviation
\end{minipage} & \begin{minipage}[b]{\linewidth}\raggedright
Status
\end{minipage} \\
\midrule\noalign{}
\endhead
\bottomrule\noalign{}
\endlastfoot
02 & FSC\_filter & 1D rectangle & Linear & 2 & 0.000000\% & 0.000000\% &
PASS \\
03 & cells & 2D rectangle & Linear & 4 & 0.000000\% & 0.000000\% &
PASS \\
04 & singlets & Polygon (6 vertices) & Linear & 6 & 0.000000\% &
0.000000\% & PASS \\
06 & FITC\_pos & 1D rectangle & Biexponential & 2 & 0.000000\% &
0.000000\% & PASS \\
07 & PE\_pos & 1D rectangle & Log & 2 & 0.000000\% & 0.000000\% & known
drift\textsuperscript{a} \\
08 & APC\_pos & 1D rectangle & Arcsinh & 2 & 0.000000\% & 0.000000\% &
PASS \\
09 & cells & 2D rectangle & Linear & 4 & 0.000000\% & 0.000000\% &
PASS \\
09 & singlets & 2D rectangle & Linear & 4 & 0.000000\% & 0.000000\% &
PASS \\
\end{longtable}

\textsuperscript{a} Test 07 uses the FlowJo logarithmic transform
(decade=6, offset=1), which exhibits coordinate drift on round-trip
between platforms. The R-side coordinates are preserved exactly in the
exported XML, but the displayed gate position in FlowJo may differ from
the original due to differences in how FlowJo and flowCore interpret log
transform parameters. See Limitations.

\textcolor{pending}{Gate coordinates are preserved with exact numerical precision (0.000\% deviation across all 7 non-drift test scenarios) for linear gate types. Deviations would only arise from floating-point rounding in the XML serialization, but the 6-decimal-place format used in the export (e.g., \texttt{60000.000000}) eliminates this source of error for integer-valued gate boundaries. However, coordinate divergence may be larger than expected in highly populated gate boundary regions; this has not been fully characterised. Ellipsoid gates (Tests 05, 12, 13) use a different XML parameterization in v10 (FlowJo foci+distance) versus flowCore (covariance matrix) and are excluded from coordinate comparison; their fidelity is assessed via population count concordance. [TODO: confirm these claims against FlowJo-side measurements.]}

\subsection{FlowJo v11 backward
compatibility}\label{flowjo-v11-backward-compatibility}

\textcolor{pending}{[TODO: Confirm all 13 exported v10 workspaces open in FlowJo v11 without errors. Record FlowJo version. If any fail, document the error and the affected gate/transform type.]}

\textcolor{pending}{[TODO: describe real-world validation with a paired v11/v10 dataset; confirm dataset availability and permissions before reporting.]}

\subsection{Standards compliance}\label{standards-compliance}

Gating-ML 2.0 compliance was assessed at two levels; full methodology
and intermediate results are provided in a dedicated vignette
(\texttt{vignette("gatingml-conformance")}). First, a structural
conformance check verified that gate elements use the correct Gating-ML
2.0 namespace URIs, element names, and required attributes. All 50
applicable checks passed across all 13 test workspaces (Table 4): 13/13
namespace declarations, 10/10 RectangleGate elements, 3/3 PolygonGate
elements, 3/3 EllipsoidGate elements, and all transformation elements
conformed to the structural requirements of the published specification
{[}4{]}.

\begin{longtable}[]{@{}llll@{}}
\caption{Table 4. Gating-ML 2.0 structural conformance results across 13
test workspaces. Each check verifies that required XML elements,
namespace URIs, and coordinate attributes conform to the published
standard (Spidlen et al.~2012).}\tabularnewline
\toprule\noalign{}
Check category & Applicable & Pass & Fail \\
\midrule\noalign{}
\endfirsthead
\toprule\noalign{}
Check category & Applicable & Pass & Fail \\
\midrule\noalign{}
\endhead
\bottomrule\noalign{}
\endlastfoot
Namespace URIs (gating, transforms, data-type) & 13 & 13 & 0 \\
Linear transform element & 13 & 13 & 0 \\
RectangleGate element & 10 & 10 & 0 \\
PolygonGate element & 3 & 3 & 0 \\
EllipsoidGate element & 3 & 3 & 0 \\
Biex transform element & 3 & 3 & 0 \\
Log transform element & 2 & 2 & 0 \\
Fasinh transform element & 3 & 3 & 0 \\
\textbf{Overall} & \textbf{50} & \textbf{50} & \textbf{0} \\
\end{longtable}

Second, strict XSD validation was attempted using the official Gating-ML
2.0 XSD files {[}4{]}. Neither CyFj11-exported WSPs nor an authentic
FlowJo v10.10.1 workspace passed strict XSD validation. In both cases,
the failure arose from the same cause: FlowJo v10's WSP format uses
proprietary transformation element names (\texttt{transforms:biex},
\texttt{transforms:linear}, \texttt{transforms:log}) that are not
defined in the Gating-ML 2.0 XSD, which expects transformation content
to be wrapped in a standard
\texttt{\textless{}transforms:transformation\textgreater{}} element.
This behavior is an inherent property of the FlowJo v10 format, not a
deficiency of CyFj11. CyFj11's output is indistinguishable from
authentic FlowJo v10 workspaces with respect to this deviation.

The MIFlowCyt standard {[}6{]} requires specific metadata for
reproducible reporting. A field-by-field audit is provided in Annex C.
CyFj11 fully preserves the MIFlowCyt fields within the scope of a gating
workspace tool: gate definitions, transformation parameters,
compensation matrices, and all FCS keyword data (including
\texttt{\$CYT}, \texttt{\$DATE}, \texttt{\$BTIM}, \texttt{\$PnN},
\texttt{\$PnS}, \texttt{\$TOT}, \texttt{\$PAR}). FCS keywords are
written to the \texttt{\textless{}Keywords\textgreater{}} section of
every sample in the v10 XML and are available in the \texttt{GatingSet}
when the user requests them via the \texttt{keywords} parameter of
\texttt{fj11\_to\_gatingset()}. Fields outside the scope of FlowJo
workspaces --- experiment purpose, antibody catalog numbers, instrument
laser configuration ---
\textcolor{pending}{may not be stored in FlowJo workspaces; see Annex C for details and verification status}.

\begin{center}\rule{0.5\linewidth}{0.5pt}\end{center}

\section{Discussion}\label{discussion}

\subsection{Practical workflow for
cytometrists}\label{practical-workflow-for-cytometrists}

CyFj11 is designed for researchers who use FlowJo v11 as their primary
gating environment and need to perform downstream analysis in R --- or
who want to export R-derived gating results back into FlowJo for visual
review or collaboration. An additional use case is researchers who
cannot use \texttt{CytoML} for importing v10 workspaces due to platform
constraints --- for example, on macOS with Apple Silicon (ARM), where
the \texttt{CytoML} Docker-based C++ components are not maintained. The
typical use case is: gate cells in FlowJo v11, import the workspace into
R with \texttt{fj11\_to\_gatingset()}, apply statistical analysis or
batch processing, and optionally return the result to FlowJo via
\texttt{export\_flowjo10\_workspace()}.

The package is not intended to replace FlowJo for gating, nor to replace
\texttt{CytoML} for users whose workflows are already based on v10
workspaces on standard platforms. It fills the specific gap left by
FlowJo's format change: the inability to read v11 workspaces in R and
the lack of a maintained, platform-independent v10 export. For v10
import, CyFj11 currently relies on \texttt{CytoML}'s FlowJo reader
(\texttt{CytoML::flowjo\_to\_gatingset()}); any limitations of that tool
--- such as partial support for certain gate types or transformations
--- also apply to the round-trip comparison workflow.

\subsection{Limitations}\label{limitations}

Several limitations should be noted. First, quadrant gates and range
gates are not exported to v10 format, because the FlowJo v10 XML schema
does not provide direct representations for these gate types; users
relying on them will need to reconstruct the gates in FlowJo after
export. Second, per-sample gate variations (tailored gates) are not
representable in v10 format; export uses the group-level gate. Because
\texttt{fj11\_to\_gatingset()} returns a list of per-sample
\texttt{GatingSet} objects rather than a single merged object, users can
inspect and work with per-sample gates in R before deciding how to merge
them; however, conversion to a single \texttt{GatingSet} via
\texttt{flowWorkspace::merge\_list\_to\_gs()} necessarily discards
sample-specific variation. Improvements to \texttt{flowCore} to natively
represent per-sample gate variations would be a valuable future
contribution. Third, gate positions in log-transformed channels undergo
a small but reproducible coordinate drift on round-trip conversion; this
is a known issue specific to the FlowJo logarithmic transform
parameterization (Test 12 in the validation suite documents this
behavior) and affects the displayed gate position, though its impact on
population counts is limited to cells near the gate boundary, which are
generally less densely populated. Fourth, the validation presented here
is based on 13 synthetic test workspaces; real-world workspaces may
contain unusual configurations not covered by these tests, though CyFj11
has been applied to large real-world datasets internally without
encountering unexpected failures. Fifth,
\textcolor{pending}{[TODO: maintenance commitment, institutional backing, community governance.]}.

\subsection{Comparison with existing
tools}\label{comparison-with-existing-tools}

\texttt{CytoML} {[}5{]} provides FlowJo v10 import via XML parsing and
export via compiled C++ components in Docker containers. CyFj11
complements this: it adds v11 import (not available in CytoML) and
provides a native-R v10 export path that does not require container
infrastructure. For users on standard hardware with v10 workspaces,
\texttt{CytoML} remains a viable choice. CyFj11's v10 export offers an
alternative for users who encounter issues with the CytoML Docker
container on current operating systems.

\texttt{openCyto} {[}3{]} provides automated gating algorithms that
operate on \texttt{GatingSet} objects; CyFj11 is fully compatible with
\texttt{openCyto} workflows --- a FlowJo v11 workspace can be imported
as a starting point for automated gating refinement.

\subsection{Future directions}\label{future-directions}

Full v11 export would require either vendor documentation of the FlowJo
v11 JSON schema or a substantially larger reverse-engineering effort
than was undertaken here. Should BD Biosciences provide official schema
documentation, implementing a v11 exporter from \texttt{GatingSet}
objects would be a natural extension of CyFj11. In the meantime, the v10
export pathway remains the practical solution for returning R analyses
to FlowJo.

Integration with spectral unmixing workflows and mass cytometry (CyTOF)
data, as well as support for emerging modalities such as imaging mass
cytometry, represent potential future directions as these data types
increasingly flow through FlowJo v11.

The package is maintained as open-source software.
\textcolor{pending}{[TODO: describe maintenance commitment, institutional backing, and community governance plan.]}

\begin{center}\rule{0.5\linewidth}{0.5pt}\end{center}

\section{Conclusion}\label{conclusion}

CyFj11 provides a practical, platform-independent solution for FlowJo
v11 / R interoperability. It enables researchers to import modern FlowJo
v11 workspaces into R with full gate and transformation support, and to
return analyzed data to FlowJo via the well-specified v10 XML format.
The pure R implementation eliminates dependencies on compiled code and
container infrastructure, ensuring compatibility across all platforms
supported by R. While the asymmetric design means that v11-specific
features such as quadrant gates and per-sample gate variations cannot be
fully round-tripped, the core gating information --- population
hierarchies, gate coordinates, compensation matrices, and transformation
parameters --- is preserved with high fidelity. CyFj11 is a practical
tool for the many cytometry laboratories that work across both the
FlowJo and R environments.

\begin{center}\rule{0.5\linewidth}{0.5pt}\end{center}

\section{References}\label{references}

\phantomsection\label{refs}
\begin{CSLReferences}{0}{0}
\bibitem[\citeproctext]{ref-hahne2009}
\CSLLeftMargin{1. }%
\CSLRightInline{Hahne F, LeMeur N, Brinkman RR, Ellis B, Haaland P,
Sarkar D, et al. flowCore: A bioconductor package for high throughput
flow cytometry. \emph{BMC bioinformatics} 2009;10:1--9.}

\bibitem[\citeproctext]{ref-finn2014}
\CSLLeftMargin{2. }%
\CSLRightInline{Finn WC, Seamer R. FlowJo v10-the latest generation of a
proven platform for flow cytometry data analysis. \emph{ASH annual
meeting abstracts} 2014;124:4400--4400.}

\bibitem[\citeproctext]{ref-finak2014}
\CSLLeftMargin{3. }%
\CSLRightInline{Finak G, Frelinger J, Jiang W, Newell EW, Ramey J, Davis
MM, et al. OpenCyto: An open source infrastructure for scalable, robust,
reproducible, and automated, end-to-end flow cytometry data analysis.
\emph{PLoS Comput Biol} 2014;10:e1003806.}

\bibitem[\citeproctext]{ref-spidlen2015}
\CSLLeftMargin{4. }%
\CSLRightInline{Spidlen J, Leif RC, Moore W, Roederer M, Brinkman RR.
Data standards: Toward consistent reporting and sharing of flow
cytometry data. \emph{Cytometry Part A} 2015;89:883--891.}

\bibitem[\citeproctext]{ref-spidlen2018}
\CSLLeftMargin{5. }%
\CSLRightInline{Spidlen J, Brinkman RR. The bioconductor flow repository
and CytoML: Infrastructure for sharing and reproducing flow cytometry
studies and results. \emph{Cytometry Part A} 2018;93:1201--1205.}

\bibitem[\citeproctext]{ref-lee2008}
\CSLLeftMargin{6. }%
\CSLRightInline{Lee JA, Spidlen J, Boyce K, Cai J, Crosbie N, Dalphin M,
et al. MIFlowCyt: The minimum information about a flow cytometry
experiment. \emph{Cytometry Part A} 2008;73:926--930.}

\end{CSLReferences}

\begin{center}\rule{0.5\linewidth}{0.5pt}\end{center}

\newpage

\section{Annex A: Technical Implementation
Details}\label{annex-a-technical-implementation-details}

\emph{This annex contains information for readers interested in the
technical architecture of CyFj11. It is not required for understanding
the package's capabilities or workflow.}

\subsection{A.1 FlowJo workspace format
comparison}\label{a.1-flowjo-workspace-format-comparison}

\textbf{FlowJo v10 (.wsp)} is a single XML file following the ISAC
Gating-ML 2.0 standard. The hierarchical structure contains
\texttt{\textless{}Workspace\textgreater{}},
\texttt{\textless{}SampleList\textgreater{}},
\texttt{\textless{}Groups\textgreater{}}, and
\texttt{\textless{}Population\textgreater{}} elements. Gates are
expressed as Gating-ML elements
(\texttt{\textless{}gating:RectangleGate\textgreater{}},
\texttt{\textless{}gating:PolygonGate\textgreater{}}, etc.) with
coordinate transformations specified in dedicated elements. The format
is human-readable, well-documented, and validated by the community.

\textbf{FlowJo v11 (.flowjo)} is a ZIP archive containing JSON files
(\texttt{analysis-*.json}, \texttt{manifest.txt}). The JSON schema uses
UUID-based cross-referencing: \texttt{populationDefinitions} (gate
geometry), \texttt{populations} (sample-specific instances with
statistics), \texttt{dataSources} (FCS file references), and
\texttt{groups} (sample collections) are separate entities linked by
UUIDs. Transformations are embedded within axis definitions rather than
stored separately. The format is proprietary and not publicly
documented.

Key structural differences that constrain the design of CyFj11:

\begin{longtable}[]{@{}
  >{\raggedright\arraybackslash}p{(\linewidth - 4\tabcolsep) * \real{0.3333}}
  >{\raggedright\arraybackslash}p{(\linewidth - 4\tabcolsep) * \real{0.3333}}
  >{\raggedright\arraybackslash}p{(\linewidth - 4\tabcolsep) * \real{0.3333}}@{}}
\toprule\noalign{}
\begin{minipage}[b]{\linewidth}\raggedright
Aspect
\end{minipage} & \begin{minipage}[b]{\linewidth}\raggedright
v10 (XML)
\end{minipage} & \begin{minipage}[b]{\linewidth}\raggedright
v11 (JSON)
\end{minipage} \\
\midrule\noalign{}
\endhead
\bottomrule\noalign{}
\endlastfoot
Format & Single XML file & ZIP with multiple JSON files \\
Standard & Gating-ML 2.0 (published) & Proprietary (undocumented) \\
Gate representation & Nested XML elements & UUID-linked JSON objects \\
Transformation storage & Separate elements & Embedded in axis
definitions \\
Per-sample variation & Not supported & desyncTable mechanism \\
Population hierarchy & Nested XML & Parent/child UUID references \\
\end{longtable}

The v11 JSON schema contains
\textcolor{pending}{[TODO: verify count manually]} top-level entity
types including: \texttt{groups}, \texttt{dataSources},
\texttt{populationDefinitions}, \texttt{populations},
\texttt{compoundParameterSets}, \texttt{compoundPopulations},
\texttt{paramsetDefinitions}, \texttt{platforms}, \texttt{cytometers},
and \texttt{analysisRoot}. Each entity maintains cross-reference
sections (\texttt{parents}, \texttt{children}, \texttt{results}) that
link to other entities by UUID. A single group entity, for example,
carries parent and child references to all other entity types, requiring
that any write operation maintain global UUID consistency across the
entire workspace graph. By contrast, the v10 XML schema uses nested
containment: a \texttt{\textless{}Population\textgreater{}} element that
contains its \texttt{\textless{}Gate\textgreater{}} is self-contained
and does not require references to be resolved elsewhere in the
document.

The complexity of constructing a valid v11 export --- requiring the
generation of a consistent UUID dependency graph across all 10 entity
types --- is the primary technical barrier to implementing v11 export.

\subsection{A.2 Gate type conversion}\label{a.2-gate-type-conversion}

\subsubsection{Rectangle gates}\label{rectangle-gates}

v11 rectangle gates define min/max boundaries per parameter. On import,
these are converted to \texttt{flowCore::rectangleGate} objects. On
export to v10, they are written as
\texttt{\textless{}gating:RectangleGate\textgreater{}} elements with
\texttt{\textless{}gating:dimension\textgreater{}} children specifying
\texttt{gating:min} and \texttt{gating:max} attributes.

1D rectangle gates (range gates in v11 terminology) are imported as
\texttt{rectangleGate} objects with one bounded and one unbounded
dimension. They are exported using the same
\texttt{\textless{}gating:RectangleGate\textgreater{}} structure.

\subsubsection{Polygon gates}\label{polygon-gates}

v11 polygon gates store vertex coordinates as arrays of
\texttt{{[}x,\ y{]}} pairs. On import, vertices are converted to a
\texttt{flowCore::polygonGate} matrix. Coordinate inverse-transformation
is applied to map from display space to data space. On export, vertices
are written as \texttt{\textless{}gating:vertex\textgreater{}} elements
within \texttt{\textless{}gating:PolygonGate\textgreater{}}.

\subsubsection{Ellipsoid gates}\label{ellipsoid-gates}

v11 ellipsoid gates are parameterized by a mean vector and covariance
matrix. On import, these are converted to
\texttt{flowCore::ellipsoidGate} objects. On export to v10, eigenvalue
decomposition is applied to derive rotation angle and semi-axis lengths
for the \texttt{\textless{}gating:EllipsoidGate\textgreater{}}
representation. This decomposition introduces minor numerical
differences relative to the original v11 gate.

\subsubsection{Boolean gates}\label{boolean-gates}

v11 Boolean gates reference other gates by UUID and specify AND/OR/NOT
combinations. On import, UUID references are resolved to population
paths and a \texttt{flowCore::booleanFilter} is constructed. On export,
Boolean logic is preserved in
\texttt{\textless{}gating:BooleanGate\textgreater{}} elements
referencing gates by population path.

\subsubsection{Quadrant gates}\label{quadrant-gates}

v11 quadrant gates define dividing lines on two axes, producing four
quadrant populations. On import, CyFj11 converts each quadrant into a
\texttt{flowCore::quadGate} object. On export to v10, quadrant gates
have no direct XML representation in the format CyFj11 generates and are
excluded with a warning.
\textcolor{pending}{[TODO: verify whether v10 WSP format itself supports quadrant gates; if so, implement export. Note: v11 JSON representation of quadrant gates should be separately verified against real FlowJo v11 files to confirm import fidelity.]}

\subsection{A.3 Transformation parameter
mapping}\label{a.3-transformation-parameter-mapping}

All transformations are extracted from v11 axis definitions and
converted to the corresponding \texttt{flowCore} transformation objects
on import. On export to v10, they are written as
\texttt{\textless{}transforms:*\textgreater{}} elements with the
appropriate parameter attributes.

\textbf{Display-space versus data-space values:} An important
distinction when working with FlowJo transformations is that FlowJo
displays gate coordinates in \emph{display space} (the transformed,
scaled axis values shown on screen), while \texttt{flowCore} and the
Gating-ML 2.0 standard store gate coordinates in \emph{data space} (raw
channel values). CyFj11 handles the inverse transformation from display
space to data space during import for gate coordinate extraction, and
applies the forward transformation during export. Users inspecting gate
coordinates directly in the v10 XML should be aware that these are in
data space and will differ from the values shown on FlowJo's axis.

\begin{longtable}[]{@{}
  >{\raggedright\arraybackslash}p{(\linewidth - 6\tabcolsep) * \real{0.2500}}
  >{\raggedright\arraybackslash}p{(\linewidth - 6\tabcolsep) * \real{0.2500}}
  >{\raggedright\arraybackslash}p{(\linewidth - 6\tabcolsep) * \real{0.2500}}
  >{\raggedright\arraybackslash}p{(\linewidth - 6\tabcolsep) * \real{0.2500}}@{}}
\toprule\noalign{}
\begin{minipage}[b]{\linewidth}\raggedright
Transformation
\end{minipage} & \begin{minipage}[b]{\linewidth}\raggedright
v11 parameter names
\end{minipage} & \begin{minipage}[b]{\linewidth}\raggedright
v10 / flowCore parameter names
\end{minipage} & \begin{minipage}[b]{\linewidth}\raggedright
Notes
\end{minipage} \\
\midrule\noalign{}
\endhead
\bottomrule\noalign{}
\endlastfoot
Biexponential & \texttt{T}, \texttt{W}, \texttt{M}, \texttt{A} &
\texttt{T}, \texttt{W}, \texttt{M}, \texttt{A} & Direct mapping \\
Logicle & \texttt{T}, \texttt{W}, \texttt{M}, \texttt{A} & \texttt{T},
\texttt{W}, \texttt{M}, \texttt{A} & Direct mapping \\
Log & \texttt{decades}, \texttt{offset} & \texttt{decades},
\texttt{offset} & Base 10 assumed \\
Arcsinh & \texttt{T}, \texttt{A} & \texttt{length}, \texttt{maxRange},
\texttt{T}, \texttt{A}, \texttt{M}, \texttt{W} & Parameter extension
required \\
Linear & \texttt{min}, \texttt{max} & \texttt{minRange},
\texttt{maxRange} & Direct mapping \\
\end{longtable}

\subsection{A.4 Package architecture}\label{a.4-package-architecture}

CyFj11 is organized into the following R source modules:

\begin{longtable}[]{@{}
  >{\raggedright\arraybackslash}p{(\linewidth - 2\tabcolsep) * \real{0.5000}}
  >{\raggedright\arraybackslash}p{(\linewidth - 2\tabcolsep) * \real{0.5000}}@{}}
\toprule\noalign{}
\begin{minipage}[b]{\linewidth}\raggedright
File
\end{minipage} & \begin{minipage}[b]{\linewidth}\raggedright
Responsibility
\end{minipage} \\
\midrule\noalign{}
\endhead
\bottomrule\noalign{}
\endlastfoot
\texttt{archive.R} & ZIP extraction, JSON parsing, manifest reading \\
\texttt{conversion.R} & Per-sample \texttt{GatingSet} list construction
from parsed workspace \\
\texttt{export-flowjo10.R} & Gating-ML 2.0 XML generation \\
\texttt{gates.R} & Gate type detection, coordinate extraction,
conversion \\
\texttt{transformations.R} & Transformation extraction and format
conversion \\
\texttt{compensation.R} & Compensation matrix extraction and validation
(see also 8 Color PBMC real-world test) \\
\texttt{populations.R} & Population hierarchy traversal and gate tree
construction \\
\texttt{file-search.R} & FCS file resolution from FlowJo URI
references \\
\texttt{helpers.R} & Utility functions, UUID validation \\
\end{longtable}

The exported functions are: \texttt{read\_flowjo11\_workspace()},
\texttt{fj11\_to\_gatingset()}, \texttt{export\_flowjo10\_workspace()},
and \texttt{pretty\_print\_flowjo()} (extracts and formats a
\texttt{.flowjo} workspace to a human-readable JSON text file, useful
for debugging). Three utility functions are also exported:
\texttt{set\_verbose(TRUE/FALSE)} enables or disables verbose progress
logging; \texttt{get\_verbose()} returns the current verbosity state;
and \texttt{\%\textbar{}\textbar{}\%} is a null-coalescing operator that
returns its left-hand argument if non-NULL, otherwise the right-hand
argument (analogous to \texttt{??} in other languages).

\subsection{A.5 Validation
methodology}\label{a.5-validation-methodology}

\subsubsection{Test workspace
construction}\label{test-workspace-construction}

Thirteen test workspaces were constructed programmatically using the
\texttt{roundtrip\_experiment.Rmd} script included with the package.
Synthetic FCS files were generated using log-normal bimodal
distributions to simulate realistic fluorescence measurements: scatter
channels (FSC-A, FSC-H, SSC-A, SSC-H) were drawn from normal
distributions (FSC-A: μ = 120,000, σ = 30,000; SSC-A: μ = 80,000, σ =
20,000), and fluorescence channels were drawn from a mixture of two
log-normal components (negative population: \textasciitilde70--80\% of
events; positive population: \textasciitilde20--30\% of events) to
produce realistic positive/negative separations. All random seeds were
fixed for reproducibility.

Gate coordinates were specified to enclose known fractions of the
simulated distributions. Transformation parameters were set to standard
FlowJo defaults: biexponential with channelRange = 4096, maxValue =
262,144, pos = 4.5, neg = 0; logicle with T = 262,144, W = 0.5, M = 4.5,
A = 0; arcsinh with T = 262,144, M = 4.5, A = 0.

\textcolor{pending}{[TODO: verify provenance of Tests 11–13 FCS files. These reference real (non-synthetic) FCS data with D-prefixed sample identifiers. Confirm the source cohort and whether these files can be included in the public deposit.]}

\textcolor{pending}{[TODO: describe real-world validation dataset for compensation matrix testing. Note: dataset under copyright; internal use only — report summary statistics only, not raw data.]}

\subsubsection{Metrics}\label{metrics}

\begin{itemize}
\tightlist
\item
  \textbf{Population count preservation:} absolute count difference and
  percentage difference between R-side \texttt{GatingSet} count and
  FlowJo-reported count after importing the exported v10 workspace.
\item
  \textbf{Gate coordinate accuracy:} mean and maximum deviation of gate
  boundary coordinates as a percentage of the axis range, computed by
  comparing coordinates in the exported v10 XML to the original gate
  definition.
\item
  \textbf{Transformation fidelity:} Kolmogorov--Smirnov D-statistic
  comparing the distribution of transformed values computed by
  \texttt{flowCore} and by FlowJo for identical parameter values and FCS
  files.
\item
  \textbf{Boolean gate accuracy:} percentage agreement in per-cell
  membership for converted Boolean gates (AND, OR, NOT), computed by
  comparing R-side and FlowJo-side population counts.
\end{itemize}

\subsubsection{Acceptance criteria}\label{acceptance-criteria}

\begin{longtable}[]{@{}ll@{}}
\toprule\noalign{}
Metric & Acceptance criterion \\
\midrule\noalign{}
\endhead
\bottomrule\noalign{}
\endlastfoot
Population count difference & ≤ 0.5\% absolute \\
Gate coordinate deviation & ≤ 0.1\% of axis range (95th percentile) \\
KS D-statistic (transformations) & \textless{} 0.05 for ≥ 95\% of
comparisons \\
Boolean membership agreement & 100\% \\
\end{longtable}

\textcolor{pending}{[TODO: run FlowJo-side validation and populate results table. Exception: Test 07 (log transform) and the log-gated populations in Test 12 are expected to fall outside the gate coordinate criterion due to documented transform drift; report the observed deviation separately for those cases.]}

\begin{center}\rule{0.5\linewidth}{0.5pt}\end{center}

\newpage

\section{Annex B: Test Scenario
Descriptions}\label{annex-b-test-scenario-descriptions}

All test scenarios use synthetic FCS data generated by the validation
suite in \texttt{roundtrip\_experiment.Rmd}. Each test directory
contains the FCS file(s), an exported v10 workspace
(\texttt{export.wsp}), and a description file (\texttt{test\_info.txt}).
\textcolor{pending}{[TODO: verify whether Tests 11–13 use real patient-derived FCS files; confirm source and data availability before public deposit.]}

\begin{longtable}[]{@{}
  >{\raggedright\arraybackslash}p{(\linewidth - 10\tabcolsep) * \real{0.1667}}
  >{\raggedright\arraybackslash}p{(\linewidth - 10\tabcolsep) * \real{0.1667}}
  >{\raggedright\arraybackslash}p{(\linewidth - 10\tabcolsep) * \real{0.1667}}
  >{\raggedright\arraybackslash}p{(\linewidth - 10\tabcolsep) * \real{0.1667}}
  >{\raggedright\arraybackslash}p{(\linewidth - 10\tabcolsep) * \real{0.1667}}
  >{\raggedright\arraybackslash}p{(\linewidth - 10\tabcolsep) * \real{0.1667}}@{}}
\toprule\noalign{}
\begin{minipage}[b]{\linewidth}\raggedright
Test
\end{minipage} & \begin{minipage}[b]{\linewidth}\raggedright
Samples
\end{minipage} & \begin{minipage}[b]{\linewidth}\raggedright
Gate hierarchy
\end{minipage} & \begin{minipage}[b]{\linewidth}\raggedright
Gate types
\end{minipage} & \begin{minipage}[b]{\linewidth}\raggedright
Transforms
\end{minipage} & \begin{minipage}[b]{\linewidth}\raggedright
Known issues
\end{minipage} \\
\midrule\noalign{}
\endhead
\bottomrule\noalign{}
\endlastfoot
01 & 1 & Root only & None & Linear & None \\
02 & 1 & 1 level & 1D rectangle (FSC-A, 60k--180k) & Linear & None \\
03 & 1 & 1 level & 2D rectangle (FSC-A, SSC-A) & Linear & None \\
04 & 1 & 1 level & Polygon on FSC-A vs FSC-H (6 vertices) & Linear &
None \\
05 & 1 & 1 level & Ellipsoid on FSC-A vs SSC-A & Linear & Minor
coordinate difference expected from eigenvalue decomposition \\
06 & 1 & 1 level & 1D rectangle on FITC-A & Biexponential
(channelRange=4096, pos=4.5, neg=0) & None \\
07 & 1 & 1 level & 1D rectangle on PE-A & Log (t=262144, m=4.5) & Gate
position drifts on round-trip; minor count difference expected \\
08 & 1 & 1 level & 1D rectangle on APC-A & Arcsinh (T=262144, M=4.5,
A=0) & None \\
09 & 1 & 2 levels: cells → singlets & 2D rectangle × 2 & Linear &
None \\
10 & 1 & 2 levels: root → \{FSC\_gate, SSC\_gate, both\} & 2× 1D
rectangle + Boolean AND & Linear & None \\
11 & 22 & 2 levels: cells → \{FITC\_pos, APC\_pos, singlets,
double\_pos, either\_pos\} & 2D rect + 2× 1D rect + polygon + 2× boolean
& Biexponential (FITC-A) + arcsinh (APC-A) + linear & Real patient
data \\
12 & 1 & 2 levels: root → \{cells, scatter\_ellipse\}; cells →
\{FITC\_hi, PE\_hi, double\_hi, either\_hi, neither\} & 2D rect +
ellipsoid + 2× range + 3× boolean (AND/OR/NOT) & Log (PE-A) +
biexponential (FITC-A) + linear & Log-transform gate drift; NOT boolean
gate tested \\
13 & 1 & 3 levels: live → \{singlets, FITC\_PE\_gate\}; singlets →
\{FITC\_pos, APC\_pos, double\_pos, not\_double\} & Polygon + 2D rect +
ellipsoid + 3× 1D rect + 2× boolean & Arcsinh (FITC-A, PE-A, APC-A) +
linear & 3-level depth; NOT boolean gate tested \\
\end{longtable}

\begin{center}\rule{0.5\linewidth}{0.5pt}\end{center}

\newpage

\section{Annex C: MIFlowCyt Compliance
Checklist}\label{annex-c-miflowcyt-compliance-checklist}

The MIFlowCyt standard {[}6{]} specifies minimum information required
for reproducible reporting of flow cytometry experiments. The table
below audits each MIFlowCyt field category against the CyFj11 pipeline:
(1) whether the information is present in a FlowJo v11 workspace, (2)
whether \texttt{fj11\_to\_gatingset()} imports it into the R
\texttt{GatingSet}, and (3) whether
\texttt{export\_flowjo10\_workspace()} writes it to the v10 XML.

Symbols: ✓ = fully preserved; \textasciitilde{} = partial / conditional;
✗ = not preserved; N/A = not applicable to software tools.
\textcolor{pending}{Items marked in red require verification.}

\begin{center}\rule{0.5\linewidth}{0.5pt}\end{center}

\subsection{C.1 Experiment Overview}\label{c.1-experiment-overview}

\begin{longtable}[]{@{}
  >{\raggedright\arraybackslash}p{(\linewidth - 8\tabcolsep) * \real{0.2000}}
  >{\raggedright\arraybackslash}p{(\linewidth - 8\tabcolsep) * \real{0.2000}}
  >{\raggedright\arraybackslash}p{(\linewidth - 8\tabcolsep) * \real{0.2000}}
  >{\raggedright\arraybackslash}p{(\linewidth - 8\tabcolsep) * \real{0.2000}}
  >{\raggedright\arraybackslash}p{(\linewidth - 8\tabcolsep) * \real{0.2000}}@{}}
\toprule\noalign{}
\begin{minipage}[b]{\linewidth}\raggedright
MIFlowCyt item
\end{minipage} & \begin{minipage}[b]{\linewidth}\raggedright
In v11 workspace?
\end{minipage} & \begin{minipage}[b]{\linewidth}\raggedright
Imported to R?
\end{minipage} & \begin{minipage}[b]{\linewidth}\raggedright
Exported to v10?
\end{minipage} & \begin{minipage}[b]{\linewidth}\raggedright
Notes
\end{minipage} \\
\midrule\noalign{}
\endhead
\bottomrule\noalign{}
\endlastfoot
Purpose / hypothesis & ✗ & ✗ & ✗ & Not stored in FlowJo workspaces or
FCS files \\
Keywords / domain terms & ✗ & ✗ & ✗ & Not stored in workspace \\
Number of experiments & N/A & N/A & N/A & Single-workspace tool \\
\end{longtable}

\begin{center}\rule{0.5\linewidth}{0.5pt}\end{center}

\subsection{C.2 Primary Specimen and Flow
Sample}\label{c.2-primary-specimen-and-flow-sample}

These fields are stored in FCS keyword fields when present; availability
depends on the instrument and acquisition software used to generate the
FCS files.

\begin{longtable}[]{@{}
  >{\raggedright\arraybackslash}p{(\linewidth - 10\tabcolsep) * \real{0.1667}}
  >{\raggedright\arraybackslash}p{(\linewidth - 10\tabcolsep) * \real{0.1667}}
  >{\raggedright\arraybackslash}p{(\linewidth - 10\tabcolsep) * \real{0.1667}}
  >{\raggedright\arraybackslash}p{(\linewidth - 10\tabcolsep) * \real{0.1667}}
  >{\raggedright\arraybackslash}p{(\linewidth - 10\tabcolsep) * \real{0.1667}}
  >{\raggedright\arraybackslash}p{(\linewidth - 10\tabcolsep) * \real{0.1667}}@{}}
\toprule\noalign{}
\begin{minipage}[b]{\linewidth}\raggedright
MIFlowCyt item
\end{minipage} & \begin{minipage}[b]{\linewidth}\raggedright
FCS keyword
\end{minipage} & \begin{minipage}[b]{\linewidth}\raggedright
In v11 workspace?
\end{minipage} & \begin{minipage}[b]{\linewidth}\raggedright
Imported to R?
\end{minipage} & \begin{minipage}[b]{\linewidth}\raggedright
Exported to v10?
\end{minipage} & \begin{minipage}[b]{\linewidth}\raggedright
Notes
\end{minipage} \\
\midrule\noalign{}
\endhead
\bottomrule\noalign{}
\endlastfoot
Organism / cell type & \texttt{\$CELLS}, \texttt{\$SRC} &
\textasciitilde{} & \textasciitilde{} & \textasciitilde{} & Available if
present in FCS keywords; extracted via \texttt{keywords} parameter \\
Sample source / tissue & \texttt{\$SRC} & \textasciitilde{} &
\textasciitilde{} & \textasciitilde{} & As above \\
Sample preparation description & --- & ✗ & ✗ & ✗ & Not stored in FCS
standard \\
Collection date and time & \texttt{\$DATE}, \texttt{\$BTIM},
\texttt{\$ETIM} & ✓ & \textasciitilde{} & ✓ & In FCS keywords; written
to v10 \texttt{\textless{}Keywords\textgreater{}} section;
\textcolor{pending}{not in `pData` by default — to be verified} \\
Operator & \texttt{\$OP} & ✓ & \textasciitilde{} & ✓ & As above \\
Institution & \texttt{\$INST} & \textasciitilde{} & \textasciitilde{} &
\textasciitilde{} & Present if recorded by acquisition software \\
Sample number / well ID & \texttt{\$SMNO}, \texttt{\$WELLID} &
\textasciitilde{} & \textasciitilde{} & \textasciitilde{} & Present if
recorded \\
\end{longtable}

\begin{center}\rule{0.5\linewidth}{0.5pt}\end{center}

\subsection{C.3 Reagents and Staining}\label{c.3-reagents-and-staining}

\begin{longtable}[]{@{}
  >{\raggedright\arraybackslash}p{(\linewidth - 10\tabcolsep) * \real{0.1667}}
  >{\raggedright\arraybackslash}p{(\linewidth - 10\tabcolsep) * \real{0.1667}}
  >{\raggedright\arraybackslash}p{(\linewidth - 10\tabcolsep) * \real{0.1667}}
  >{\raggedright\arraybackslash}p{(\linewidth - 10\tabcolsep) * \real{0.1667}}
  >{\raggedright\arraybackslash}p{(\linewidth - 10\tabcolsep) * \real{0.1667}}
  >{\raggedright\arraybackslash}p{(\linewidth - 10\tabcolsep) * \real{0.1667}}@{}}
\toprule\noalign{}
\begin{minipage}[b]{\linewidth}\raggedright
MIFlowCyt item
\end{minipage} & \begin{minipage}[b]{\linewidth}\raggedright
FCS keyword
\end{minipage} & \begin{minipage}[b]{\linewidth}\raggedright
In v11 workspace?
\end{minipage} & \begin{minipage}[b]{\linewidth}\raggedright
Imported to R?
\end{minipage} & \begin{minipage}[b]{\linewidth}\raggedright
Exported to v10?
\end{minipage} & \begin{minipage}[b]{\linewidth}\raggedright
Notes
\end{minipage} \\
\midrule\noalign{}
\endhead
\bottomrule\noalign{}
\endlastfoot
Fluorochrome names & \texttt{\$PnS} & ✓ & ✓ & ✓ & Channel stain
descriptions; in parameter keywords \\
Channel names (detector) & \texttt{\$PnN} & ✓ & ✓ & ✓ & Required FCS
keyword; always present \\
Antibody clone / catalog & --- & ✗ & ✗ & ✗ & Not a standard FCS field \\
Antibody concentration & --- & ✗ & ✗ & ✗ & Not a standard FCS field \\
Viability dye & --- & ✗ & ✗ & ✗ & Convention only; not standardized \\
\end{longtable}

\begin{center}\rule{0.5\linewidth}{0.5pt}\end{center}

\subsection{C.4 Instrument}\label{c.4-instrument}

\begin{longtable}[]{@{}
  >{\raggedright\arraybackslash}p{(\linewidth - 10\tabcolsep) * \real{0.1667}}
  >{\raggedright\arraybackslash}p{(\linewidth - 10\tabcolsep) * \real{0.1667}}
  >{\raggedright\arraybackslash}p{(\linewidth - 10\tabcolsep) * \real{0.1667}}
  >{\raggedright\arraybackslash}p{(\linewidth - 10\tabcolsep) * \real{0.1667}}
  >{\raggedright\arraybackslash}p{(\linewidth - 10\tabcolsep) * \real{0.1667}}
  >{\raggedright\arraybackslash}p{(\linewidth - 10\tabcolsep) * \real{0.1667}}@{}}
\toprule\noalign{}
\begin{minipage}[b]{\linewidth}\raggedright
MIFlowCyt item
\end{minipage} & \begin{minipage}[b]{\linewidth}\raggedright
FCS keyword
\end{minipage} & \begin{minipage}[b]{\linewidth}\raggedright
In v11 workspace?
\end{minipage} & \begin{minipage}[b]{\linewidth}\raggedright
Imported to R?
\end{minipage} & \begin{minipage}[b]{\linewidth}\raggedright
Exported to v10?
\end{minipage} & \begin{minipage}[b]{\linewidth}\raggedright
Notes
\end{minipage} \\
\midrule\noalign{}
\endhead
\bottomrule\noalign{}
\endlastfoot
Cytometer manufacturer and model & \texttt{\$CYT} & ✓ &
\textasciitilde{} & ✓ & In FCS keywords; written to v10 XML
\texttt{\textless{}Keywords\textgreater{}};
\textcolor{pending}{not in `pData` by default — to be verified} \\
Excitation laser wavelengths & --- & ✗ & ✗ & ✗ & Not in FCS standard;
instrument-specific \\
Detector/filter configuration & \texttt{\$PnF} (non-standard) &
\textasciitilde{} & \textasciitilde{} & \textasciitilde{} &
Instrument-specific; not always recorded \\
Amplifier gain / voltage & \texttt{\$PnV}, \texttt{\$PnG} &
\textasciitilde{} & \textasciitilde{} & \textasciitilde{} &
Instrument-specific; present if recorded \\
\end{longtable}

\begin{center}\rule{0.5\linewidth}{0.5pt}\end{center}

\subsection{C.5 Acquisition}\label{c.5-acquisition}

\begin{longtable}[]{@{}
  >{\raggedright\arraybackslash}p{(\linewidth - 10\tabcolsep) * \real{0.1667}}
  >{\raggedright\arraybackslash}p{(\linewidth - 10\tabcolsep) * \real{0.1667}}
  >{\raggedright\arraybackslash}p{(\linewidth - 10\tabcolsep) * \real{0.1667}}
  >{\raggedright\arraybackslash}p{(\linewidth - 10\tabcolsep) * \real{0.1667}}
  >{\raggedright\arraybackslash}p{(\linewidth - 10\tabcolsep) * \real{0.1667}}
  >{\raggedright\arraybackslash}p{(\linewidth - 10\tabcolsep) * \real{0.1667}}@{}}
\toprule\noalign{}
\begin{minipage}[b]{\linewidth}\raggedright
MIFlowCyt item
\end{minipage} & \begin{minipage}[b]{\linewidth}\raggedright
FCS keyword
\end{minipage} & \begin{minipage}[b]{\linewidth}\raggedright
In v11 workspace?
\end{minipage} & \begin{minipage}[b]{\linewidth}\raggedright
Imported to R?
\end{minipage} & \begin{minipage}[b]{\linewidth}\raggedright
Exported to v10?
\end{minipage} & \begin{minipage}[b]{\linewidth}\raggedright
Notes
\end{minipage} \\
\midrule\noalign{}
\endhead
\bottomrule\noalign{}
\endlastfoot
Trigger / threshold channel & --- & ✗ & ✗ & ✗ & Not standardized in
FCS \\
Threshold value & --- & ✗ & ✗ & ✗ & Not standardized in FCS \\
Collection volume / events & \texttt{\$TOT} & ✓ & ✓ & ✓ & Total event
count; used in sample naming \\
Number of parameters & \texttt{\$PAR} & ✓ & ✓ & ✓ & Required FCS
keyword \\
FCS file format version & \texttt{FCSversion} & ✓ & ✓ & ✓ & Written as
keyword \\
Data type & \texttt{\$DATATYPE} & ✓ & \textasciitilde{} & ✓ & In FCS
keywords \\
Byte order & \texttt{\$BYTEORD} & ✓ & \textasciitilde{} & ✓ & In FCS
keywords \\
\end{longtable}

\begin{center}\rule{0.5\linewidth}{0.5pt}\end{center}

\subsection{C.6 Data Analysis}\label{c.6-data-analysis}

\begin{longtable}[]{@{}
  >{\raggedright\arraybackslash}p{(\linewidth - 10\tabcolsep) * \real{0.1667}}
  >{\raggedright\arraybackslash}p{(\linewidth - 10\tabcolsep) * \real{0.1667}}
  >{\raggedright\arraybackslash}p{(\linewidth - 10\tabcolsep) * \real{0.1667}}
  >{\raggedright\arraybackslash}p{(\linewidth - 10\tabcolsep) * \real{0.1667}}
  >{\raggedright\arraybackslash}p{(\linewidth - 10\tabcolsep) * \real{0.1667}}
  >{\raggedright\arraybackslash}p{(\linewidth - 10\tabcolsep) * \real{0.1667}}@{}}
\toprule\noalign{}
\begin{minipage}[b]{\linewidth}\raggedright
MIFlowCyt item
\end{minipage} & \begin{minipage}[b]{\linewidth}\raggedright
Source
\end{minipage} & \begin{minipage}[b]{\linewidth}\raggedright
In v11 workspace?
\end{minipage} & \begin{minipage}[b]{\linewidth}\raggedright
Imported to R?
\end{minipage} & \begin{minipage}[b]{\linewidth}\raggedright
Exported to v10?
\end{minipage} & \begin{minipage}[b]{\linewidth}\raggedright
Notes
\end{minipage} \\
\midrule\noalign{}
\endhead
\bottomrule\noalign{}
\endlastfoot
Analysis software name & Workspace JSON & ✓ & ✗ & ✗ & FlowJo version in
\texttt{schemaVersion}; not propagated to GatingSet \\
Analysis software version & Workspace JSON & ✓ & ✗ & ✗ & As above \\
Gate names & populationDefinitions & ✓ & ✓ & ✓ & Fully extracted into
\texttt{GatingHierarchy} \\
Gate coordinates & gateDefinition & ✓ & ✓ & ✓ & All gate types including
per-sample variations \\
Gate hierarchy & parent/child UUIDs & ✓ & ✓ & ✓ & Full population tree
preserved \\
Transformation type & axis parameterSpec & ✓ & ✓ & ✓ & biexp, logicle,
log, arcsinh, linear \\
Transformation parameters & axis parameterSpec & ✓ & ✓ & ✓ & T, W, M, A
parameters for each channel \\
Compensation matrix & platforms.spilloverMatrix & ✓ & ✓ & ✓ & Extracted
from multiple possible sources \\
Compensation controls & Compensation group & ✓ & \textasciitilde{} &
\textasciitilde{} & Group identified; FCS files resolved \\
Population statistics & populations results & ✓ & ✓ & ✓ & Counts,
frequencies after gate execution \\
Per-sample gate variations & desyncTable & ✓ & ✓ & \textasciitilde{} &
Imported; merged to group gate on export \\
\end{longtable}

\begin{center}\rule{0.5\linewidth}{0.5pt}\end{center}

\subsection{C.7 Summary}\label{c.7-summary}

\begin{longtable}[]{@{}
  >{\raggedright\arraybackslash}p{(\linewidth - 6\tabcolsep) * \real{0.2500}}
  >{\raggedright\arraybackslash}p{(\linewidth - 6\tabcolsep) * \real{0.2500}}
  >{\raggedright\arraybackslash}p{(\linewidth - 6\tabcolsep) * \real{0.2500}}
  >{\raggedright\arraybackslash}p{(\linewidth - 6\tabcolsep) * \real{0.2500}}@{}}
\toprule\noalign{}
\begin{minipage}[b]{\linewidth}\raggedright
Category
\end{minipage} & \begin{minipage}[b]{\linewidth}\raggedright
Import fidelity
\end{minipage} & \begin{minipage}[b]{\linewidth}\raggedright
Export fidelity
\end{minipage} & \begin{minipage}[b]{\linewidth}\raggedright
Key limitation
\end{minipage} \\
\midrule\noalign{}
\endhead
\bottomrule\noalign{}
\endlastfoot
Experiment overview & ✗ & ✗ & Not stored in FlowJo workspaces \\
Sample metadata & \textasciitilde{} & \textasciitilde{} & FCS keyword
fields are preserved but not surfaced to \texttt{pData} by default; use
\texttt{keywords} parameter \\
Reagents / staining & \textasciitilde{} & \textasciitilde{} &
\texttt{\$PnS} (stain names) preserved; antibody details not in FCS
standard \\
Instrument & \textasciitilde{} & \textasciitilde{} & \texttt{\$CYT} and
parameter keywords preserved in XML; not in \texttt{pData} by default \\
Acquisition & ✓ & ✓ & Core fields (\texttt{\$TOT}, \texttt{\$PAR},
\texttt{\$PnN}) always preserved \\
Data analysis (gates) & ✓ & ✓ & Complete: hierarchy, coordinates, types,
transforms \\
Compensation & ✓ & ✓ & Full matrix from multiple extraction paths \\
Software provenance & \textasciitilde{} & ✗ & FlowJo version available
in v11 JSON; not written to v10 XML \\
\end{longtable}

\textbf{Overall assessment:} CyFj11 fully preserves the MIFlowCyt fields
that are within the scope of a gating workspace tool --- gate
definitions, transformations, compensation matrices, and FCS keyword
data.
\textcolor{pending}{Fields outside the FCS/workspace scope (experiment purpose, antibody catalogs, instrument laser configuration) are generally not stored in FlowJo workspaces, but this should be verified against real-world workspace files (e.g., the 8 Color PBMC dataset) to confirm.}

\end{document}
